% GENERATED FILE. Edit canonical lessons, then run npm run generate:books.
% canonical-source-hash: fnv1a64:ead0f238ea2689a8
% canonical-lessons: SA-C265-meal-and-restaurant-recall, SA-C265-dadatu, SA-C265-give-me-an-item, SA-C265-polite-order-contrast, SA-C265-mixed-order-practice, SA-C265-restaurant-order-retrieval

\chapter{A gentle order at a meal}
\label{ch:sa-gentle-meal-order}
\hlchaptermodality{265}

\begin{chapteropening}
\textbf{By the end of this chapter:} \emph{I can politely ask for a familiar food or drink in a restaurant.}

Complete chapter 265 by independently choosing a familiar item and making a polite request.
\end{chapteropening}

\section[Practice]{A meal in a familiar place}

\label{lesson:SA-C265-meal-and-restaurant-recall}



\begin{quote}
Say “armour,” then contrast “my” with “our,” and finish with “I will do.” (\emph{Kavacam; mama, asmākam; kariṣyāmi.})
\end{quote}

\begin{grammarlens}[title={Grammar Lens: bring the meal words back}]
Nothing here is new. \textbf{\sk{भोजनालयः}} (\emph{bhojanālayaḥ}) is a restaurant; \textbf{\sk{भोजनम्}} (\emph{bhojanam}) is a meal or food. You already know \textbf{\sk{खादति}} (\emph{khādati}), “eats,” along with \textbf{\sk{अन्नम्}} (\emph{annam}), food, \textbf{\sk{पानीयम्}} (\emph{pānīyam}), water, and \textbf{\sk{क्षीरम्}} (\emph{kṣīram}), milk.
\end{grammarlens}

\subsection*{Your turn}
\begin{itemize}
\raggedright
  \item \emph{Say it:} “restaurant” — \emph{bhojanālayaḥ}
  \item \emph{Say it:} “meal” — \emph{bhojanam}
  \item \emph{Say it:} “eats food” — \emph{annaṃ khādati}
  \item \emph{Contrast:} water, then milk — \emph{pānīyam; kṣīram}
\end{itemize}

\begin{tcolorbox}[breakable,colback=teal!4,colframe=teal!35!black,title={Before you move on}]
Say “restaurant” and “meal.” (\textbf{\sk{भोजनालयः}; \sk{भोजनम्}}.)
\end{tcolorbox}
\section[\texorpdfstring{\sk{ददातु} (dadātu)}{ददातु (dadātu)}]{\sk{ददातु} (dadātu) — please give}

\label{lesson:SA-C265-dadatu}



\begin{quote}
Say “flag,” then “me” as an object, then “I will read.” (\emph{Dhvajaḥ; mām; paṭhiṣyāmi.})
\end{quote}

\begin{grammarlens}[title={Grammar Lens: one useful request form}]
You know \textbf{\sk{ददाति}} (\emph{dadāti}), “he or she gives.” For a gentle request, use the complete form \textbf{\sk{ददातु}} (\emph{dadātu}): “let him or her give,” or in this exchange, “please give.”

\begin{quote}
\textbf{\sk{पानीयम्} \sk{ददातु।}}
\end{quote}

\begin{quote}
\emph{Pānīyaṃ dadātu.}
\end{quote}

\begin{quote}
Please give water.
\end{quote}

Here \textbf{\sk{ददातु}} is the \textbf{third-person singular imperative}. Whitney, \emph{Sanskrit Grammar} §§617 and 668–669, records the present imperative and this reduplicating root. Keep the whole form together. This is \textbf{not a complete imperative paradigm}: other persons, numbers, and roots wait for later.
\end{grammarlens}

\subsection*{Your turn}
\begin{itemize}
\raggedright
  \item \emph{Contrast:} gives; please give — \emph{dadāti; dadātu}
  \item \emph{Say it:} “Please give water” — \emph{pānīyaṃ dadātu}
  \item \emph{Say it:} the request form once more — \emph{dadātu}
\end{itemize}

\begin{tcolorbox}[breakable,colback=teal!4,colframe=teal!35!black,title={Before you move on}]
Change \textbf{\sk{ददाति}}, “gives,” into the one taught request form. (\textbf{\sk{ददातु}}.)
\end{tcolorbox}
\section[(give me a familiar item)]{Give me one familiar item}

\label{lesson:SA-C265-give-me-an-item}



\begin{quote}
Say “drum,” then “you” as one object, then make one plan for tomorrow. (\emph{Paṭahaḥ; tvām; śvaḥ} plus a taught first-person future form.)
\end{quote}

\begin{grammarlens}[title={Grammar Lens: receiver, item, request}]
Build the request in three small pieces:

\begin{quote}
\textbf{\sk{मह्यम्}} + familiar item + \textbf{\sk{ददातु}}
\end{quote}

\begin{quote}
to me + item + please give
\end{quote}

So \textbf{\sk{मह्यम्} \sk{पानीयं} \sk{ददातु}} (\emph{mahyaṃ pānīyaṃ dadātu}) asks, “Please give me water.” The item changes; the frame stays still.
\end{grammarlens}

\subsection*{Your turn}
Say these aloud:

\begin{itemize}
\raggedright
  \item “Give me food” — \emph{mahyaṃ annaṃ dadātu}
  \item “Give me water” — \emph{mahyaṃ pānīyaṃ dadātu}
  \item “Give me milk” — \emph{mahyaṃ kṣīraṃ dadātu}
\end{itemize}

\begin{tcolorbox}[breakable,colback=teal!4,colframe=teal!35!black,title={Before you move on}]
What are the three pieces? (Receiver \textbf{\sk{मह्यम्}}, a familiar item, then \textbf{\sk{ददातु}}.)
\end{tcolorbox}
\section[Practice]{Name the meal or make the order}

\label{lesson:SA-C265-polite-order-contrast}



\begin{quote}
Say “treasury,” then explain the choice between \textbf{\sk{माम्}} and \textbf{\sk{त्वाम्}}: is the one person used as object the speaker or the addressee?
\end{quote}

\begin{grammarlens}[title={Grammar Lens: name or request}]
\textbf{\sk{भोजनम्}} names a meal. It does not itself order anything. Put the courtesy word before the frame when you make the request:

\begin{quote}
\textbf{\sk{कृपया} \sk{मह्यम्} \sk{अन्नं} \sk{ददातु।}}
\end{quote}

\begin{quote}
\emph{Kṛpayā mahyam annaṃ dadātu.}
\end{quote}

\begin{quote}
Please give me food.
\end{quote}
\end{grammarlens}

\subsection*{Your turn}
Say these aloud:

\begin{itemize}
\raggedright
  \item name a meal — \emph{bhojanam}
  \item “Please give me water” — \emph{kṛpayā mahyaṃ pānīyaṃ dadātu}
  \item “Please give me milk” — \emph{kṛpayā mahyaṃ kṣīraṃ dadātu}
\end{itemize}

\begin{tcolorbox}[breakable,colback=teal!4,colframe=teal!35!black,title={Before you move on}]
Add the courtesy word to one complete request of your choice.
\end{tcolorbox}
\section[Practice]{Choose and order}

\label{lesson:SA-C265-mixed-order-practice}



\begin{quote}
Say “treasure,” then say the independent object form for “us two.” (\emph{Nidhiḥ; āvām.})
\end{quote}

\subsection*{Your turn — independent choice}
No Sanskrit model is shown before you answer. Build each request from meaning.

Say these aloud:

\begin{itemize}
\raggedright
  \item In a restaurant, politely ask for water.
  \item Politely ask for milk.
  \item Politely ask for food.
  \item Name the place, name the meal, then say that someone eats food.
\end{itemize}

\begin{grammarlens}[title={What you've built}]
Possible answers are \textbf{\sk{कृपया} \sk{मह्यं} \sk{पानीयं} \sk{ददातु}}, \textbf{\sk{कृपया} \sk{मह्यं} \sk{क्षीरं} \sk{ददातु}}, and \textbf{\sk{कृपया} \sk{मह्यम्} \sk{अन्नं} \sk{ददातु}}. The place is \textbf{\sk{भोजनालयः}}; the meal is \textbf{\sk{भोजनम्}}; “eats food” is \textbf{\sk{अन्नं} \sk{खादति}}.
\end{grammarlens}

\begin{tcolorbox}[breakable,colback=teal!4,colframe=teal!35!black,title={Before you move on}]
Choose a different item from your last answer and order it again.
\end{tcolorbox}
\section[Practice]{Make the order yourself}

\label{lesson:SA-C265-restaurant-order-retrieval}



\begin{quote}
Say “rupee,” then say the independent object form for “you two.” (\emph{Rūpyakam; yuvām.})
\end{quote}

\subsection*{Your turn — final retrieval}
No Sanskrit model is shown before you answer. Imagine you are in a \textbf{\sk{भोजनालयः}}. Choose \textbf{\sk{अन्नम्}}, \textbf{\sk{पानीयम्}}, or \textbf{\sk{क्षीरम्}}, and make the complete polite request from memory.

Now choose a second item and make a different request.

Finally, name the meal and say that someone eats food.

\begin{grammarlens}[title={What you've built}]
Your request should contain \textbf{\sk{कृपया}}, \textbf{\sk{मह्यम्}}, the chosen item, and \textbf{\sk{ददातु}}. For example, after answering: \textbf{\sk{कृपया} \sk{मह्यं} \sk{जलं} \sk{ददातु}} is the mock-paper water request. \textbf{\sk{भोजनम्}; \sk{अन्नं} \sk{खादति}} completes the last prompt.
\end{grammarlens}

\begin{grammarlens}[title={Grammar Lens: the boundary}]
This chapter gives you one useful request with familiar items. It is \textbf{not a complete imperative paradigm}, a complete restaurant vocabulary, or evidence that every possible meal exchange is covered.
\end{grammarlens}

\begin{tcolorbox}[breakable,colback=teal!4,colframe=teal!35!black,title={Before you move on}]
Without looking back, make one final polite order for a familiar food or drink.
\end{tcolorbox}
